\begin{scanpage}{0016}{xii}
Let us note the following feature of the method discussed in the book. A physical representation of the nonuniform part of the current is widely used in the book, but nowhere are its explicit mathematical expressions cited. This part of the current is generally not expressed in terms of well-known functions. Obviously a direct integration of the currents when calculating the fringing field is able to lead only to very complicated and immense equations. Therefore, we will find the fringing field created by the nonuniform part of the current on the basis of indirect considerations without direct integration of it (see especially Chapters I--IV).

The method by which the diffraction problems are solved in this book may be briefly summarized as follows. We will seek an approximate solution of the diffraction problem for a certain body by first having studied diffraction by its separate geometric elements. For example, for a finite cylinder such elements are: the lateral surface as part of an infinite cylindrical surface, each base as part of a plane, each section of the base rim as the edge of a wedge (the curvature of the rim in the first approximation may be neglected). Having studied the diffraction by the separate elements of the body, we will obtain a representation of the nonuniform part of the current and of the field which is radiated by it. Then secondary, tertiary, etc.\ diffraction is studied --- that is, the diffraction interaction of the various elements of the body is taken into account.

This method appeals to physical considerations, not only when formulating the problem but also in its solution process, and in this way differs from the methods of the mathematical theory of diffraction. Therefore, such a method may be referred to as the \uline{physical} theory of diffraction.

A whole series of other diffraction studies which appeared in the last five to ten years also are able to relate to the physical theory of diffraction. The first work which contained the idea of the physical theory of diffraction is evidently the paper of Schwarzschild [15] which was published at the beginning of this century and was devoted to diffraction by a slit.
\end{scanpage}
